% ===== sec/06_implications.tex =====
\section{Design Implications and Framework}
\label{sec:implications}

The trace-language account yields three design principles for
practical agent architectures and a reusable verification framework.

\subsection{Implication 1: Separate generation from validation}

An LLM that both proposes and judges its own output conflates the
generator and verifier roles.  When a single Type-0 machine serves
both functions, the verifier cost is bounded only by the generator's
complexity, and the agent cannot certify its own trace in any
complexity-theoretic sense.  The architectural fix is a strict
separation:

\[
\mathcal{G}\;(\text{Type-0}) \quad\longrightarrow\quad
\mathcal{V}\;(\text{Type-3})
\]

where $\mathcal{V}$ is a DFA with $O(n)$ decision cost regardless of
$\mathcal{G}$'s internal complexity.  This separation is not merely
good engineering: it is a consequence of
Theorem~\ref{thm:closure}, which guarantees that the composition
remains in $\mathcal{G}$'s class.  Every agent system that uses an
LLM as its own judge implicitly places the \emph{verifier} in
Type-0, forfeiting the cost bound of
Corollary~\ref{cor:budget}.

\subsection{Implication 2: Keep validators lower in the hierarchy}

When multiple validation stages are composed, each should be no
higher in the Chomsky hierarchy than the stage it constrains.  A
Type-1 (LBA) validator is decidable but exponential; a Type-3 (DFA)
validator is decidable in linear time.  The ablation results
(Section~\ref{sec:results}) demonstrate this empirically: replacing
the DFA verifier with an LLM judge ($\mathcal{G}\,\Vert\,\mathcal{V}_{\mathrm{LLM}}$)
raised completion to $100\%$ but did \emph{not} improve accuracy
and increased verifier latency from $68$\,ms to $14.7$\,s---a
$216\times$ slowdown.  The DFA verifier's class bound is not a
theoretical nicety; it translates directly to runtime cost.

\subsection{Implication 3: Agent roles emerge from traces}

Sub-agent labels---planner, reviewer, researcher, trainer---are
projections of the validated global trace onto sub-alphabets
(Theorem~\ref{thm:derived}).  They are not primitives that must be
pre-engineered.  This means the framework can discover role
structure: given only $\mathcal{G}$, $\mathcal{V}$, and the start
alphabet $\Sigma_0$, the recursive enumeration
(Algorithm~\ref{alg:enum}) produced a thirteenth projected column
(\texttt{kaggle\_trainer}) that was not in the initial design.
The role emerged from the verifier's pressure, not from an explicit
architectural decision.  In this sense the trace-language account
reverses the conventional workflow: instead of designing roles and
then verifying them, one designs the verifier and lets the roles
fall out as projections of the resulting fixed point.

\begin{figure}[ht]
\centering
\vspace{2pt}
\framebox[\columnwidth]{\centering\strut\textbf{Trace Recorder}}
\vspace{4pt}

\noindent\hfil$\downarrow$\hfil
\vspace{4pt}

\framebox[\columnwidth]{\centering\strut\textbf{DFA Verifier}}
\vspace{4pt}

\noindent\hfil$\downarrow$\hfil
\vspace{4pt}

\framebox[\columnwidth]{\centering\strut\textbf{Projection Engine}}
\vspace{4pt}

\noindent\hfil$\downarrow$\hfil
\vspace{4pt}

\framebox[\columnwidth]{\centering\strut\textbf{Agent Analysis}}
\vspace{2pt}
\caption{Trace-language verification framework.  A practitioner
records traces from any agent system, filters them through a
task-specific DFA verifier, projects onto sub-alphabets, and
analyses the resulting language-theoretic class signature.}
\label{fig:framework}
\end{figure}

\subsection{Reusable framework}

Figure~\ref{fig:framework} distils the methodology into four
stages that apply to any agent system, not only the
\textsc{Spaceship Titanic} instantiation:

\paragraph{Trace Recorder.}
Instrument the agent to emit a log of operation symbols---tool
calls, API invocations, state transitions.  The output is a CSV
whose columns are sub-agent traces.

\paragraph{DFA Verifier.}
Compile task-specific unit tests and schema constraints into an
Aho--Corasick keyword automaton
(Proposition~\ref{prop:verifier-dfa}).  The verifier accepts or
rejects each candidate trace in $O(n)$ time.

\paragraph{Projection Engine.}
Decompose the validated global trace into sub-alphabets
corresponding to logical roles.  Each projection is an empirical
sample of a sub-agent's trace language.

\paragraph{Agent Analysis.}
Assign Chomsky-class hypotheses to each projected trace
(Table~\ref{tab:class-signature}).  The verifier's own recogniser is
concrete (Type-3 DFA); all other classes are modelling hypotheses
that guide the search for more precise recognisers.

The framework is agnostic to the agent implementation language,
LLM provider, and task domain.  The only task-specific artefacts
are the start alphabet $\Sigma_0$ and the keyword set $K$ of the
DFA verifier.
